\ssscp{\tsc{Havu-Dhao} phnology}{03-03-01-02}
Like the \tsc{Flores} languages
Havu and Dhao have schwa, but schwa in these two languages
has different behaviour compared to
languages of Flores (\srf{03-02-01-02}).
The three main ways in which it differs are:

\begin{itemize}
	\item Schwa has normal duration and full syllable weight.
	\item Schwa is permitted in vowel sequences followed by a high vowel.
	\item Schwa regularly bears stress when penultimate.
\end{itemize}

Like the languages of Flores, schwa in Havu and Dhao
triggers phonetic lengthening of the following consonant
and consonant lengthening is not triggered by any other vowel.
Lengthening of consonants after schwa is a common feature of Austronesian
\citep{Smith2023}.
Minimal pairs from Havu showing the
behaviour of schwa are given in \trf{HavSchOthVowCon}
(See \srf{03-03-02-HD} for the development of schwa in \tsc{Havu-Dhao}).

\begin{table}
	\centering\tcp{Havu schwa in contrast with other vowels}{HavSchOthVowCon}
	\st{6pt}\begin{tabular}{lll|lll}\lsptoprule
	Havu	&	Phonetic&	Gloss										&	Havu		&	Phonetic	&	Gloss			\\	\midrule
	ɓəla	&	[ˈɓəlːɐ]	&	cloth									& ɓela		&	[ˈɓelɐ]		&	lightning			\\
	əle		&	[ˈʔəlːe]	&	finish, complete			&	ele			&	[ˈʔele]		&	disappear	\\
	həɓa	&	[ˈhəɓːɐ]	&	mouth, opening				&	haɓa		&	[ˈhaɓɐ]		&	work			\\
	ʄəge	&	[ˈʄəgːe]	&	work, job,						&	ʄege		&	[ˈʄege]		&	snarl, rebuke,	\\
				&						&	responsibility				&					&						&	mad at	\\
%	məhu	&	[ˈməhːu]	&	exit, come out from		&	maho		&	[ˈmaho]		&	enter			\\
	vəla	&	[ˈvəlːɐ]	&	spread open 					&	vela		&	[ˈvelɐ]		&	machete		\\
				&						&	and place on					&					&						&						\\
	ɗəi		&	[ˈɗəi]		&	will, desire					&	ɗai			&	[ˈɗai]		&	arrive		\\
	əi		&	[ˈʔəi]		&	water									&	ai			&	[ˈʔai]		&	fire			\\
%												&	vala	&	[ˈvalɐ]	&	‘other
%												&	ɓala	&	[ˈɓalɐ]	&	‘reciprocate’
		\lspbottomrule
	\end{tabular}
\end{table}

Schwa has multiple sources in \tsc{Havu-Dhao},
this includes retention of *ə = \it{ə} in penultimate syllables,
as well as syllable-final *ə and *a, in which case the
schwa is associated with vowel-metathesis.
See \srf{03-03-02-HD} for details and examples.

Regarding consonants, Havu and Dhao have 
an imploded palatal plosive /ʄ/ ‹j{\Q}›
and Dhao has a slightly retroflex lightly
articulated alveolar affricate /ɖ͡ʐ/ ‹dh›,
as well as a light bilabial affricate /b͡β/ ‹bh›
\citep[256]{Grimes2010a}.
